\documentclass[../../../include/open-logic-section]{subfiles}

\begin{document}

\olfileid{sth}{replacement}{absinf}
\olsection{Panggantian lan ``Tanpa Wates Mutlak''}

Saiki kita ninggalake \limofsize{}, banjur
nliti kulawarga upaya (intrinsik) liyane
kanggo menehi dhasar marang Panggantian,
kang nganggep kanthi tenanan gagasan yen
himpunan kabentuk tataran demi tataran.

Nalika pisanan nggambarake proses iteratif,
kita ngusulake prinsip-prinsip kang nerangake
apa kang kelakon ing saben tataran: yaiku
\stageshier, \stagesord, lan \stagesacc.
Sabanjure kita nambah prinsip-prinsip babagan
cacahing tataran: \stagessucc{} ngandhakake
yen proses pambentukan himpunan ora tau
rampung, lan \stagesinf{} ngandhakake yen
proses iki tekan tataran tanpa wates kang
kapisan.

Lumrah yen rong prinsip pungkasan iku
dianggep ndherek saka prinsip kang luwih amba,
kayata:
\begin{enumerate}
	\item[] \stagesinex. Ana tataran kang tanpa
	wates kanthi mutlak; hierarki iku sa-dhuwur
	kang bisa digayuh.
\end{enumerate}
Cetha yen iki prinsip ora formal. Sanadyan
ora langsung \emph{ndherek} saka konsepsi
himpunan kang kumulatif lan iteratif, prinsip
iki katon \emph{selaras} karo konsepsi mau.
Paling ora, beda karo \limofsize, prinsip iki
njaga gagasan yen himpunan kabentuk tataran
demi tataran.

Saiki pangarep-arepe yaiku nggunakake
\stagesinex{} kanggo menehi dhasar marang
Panggantian. Ayo dideleng carane.

Ing \olref[ordinals][idea]{sec}, kita weruh
yen gampang mbangun urutan becik kang miturut
intuisi mesthine isomorfis karo
$\omega+\omega$. Kanthi tembung liya, kita
bisa mbayangake proses iteratif tataran demi
tataran kang tipe urutane, miturut intuisi,
$\omega+\omega$. Dadi yen wis nampa
\stagesinex, kita mesthine uga nampa anane
paling ora tataran kaping-$\omega+\omega$
ing hierarki, yaiku $V_{\omega+\omega}$,
amarga hierarki iku mesthi \emph{bisa}
terus tekan kono.

Gagasan iki bisa diumumake mangkene: kanggo
saben urutan becik, proses pambangunan hierarki
iteratif kudu tekan paling ora sasuwene urutan
becik iku. Iki bisa dijamin yen
\olref[ordinals][ordtype]{thmOrdinalRepresentation}
dianggep minangka \emph{aksioma}. Mula saben
urutan becik bakal isomorfis karo ordinal von
Neumann. Amarga rank saben ordinal von Neumann
padha karo ordinal iku dhewe,
\olref[ordinals][ordtype]{thmOrdinalRepresentation}
bakal ngandhakake yen kapan wae kita bisa
njlentrehake urutan becik ing teori himpunan,
proses pambentukan himpunan kang iteratif kudu
ngluwihi urutan becik mau.

Gagasan iki katon minangka akibat saka
\stagesinex. Nanging yen tujuwane njupuk
Panggantian saka gagasan iki, ana alangan
teknis kang prasaja: Panggantian luwih kuwat
tinimbang
\olref[ordinals][ordtype]{thmOrdinalRepresentation}.
(Iki dicathet dening \citet[\S13.2]{Potter2004};
kita bakal mbuktekake ing
\olref[sth][replacement][finiteaxiomatizability]{sec}.)

Dadi yen \stagesinex{} arep dimangerteni
kanthi cara kang ngasilake Panggantian,
prinsip iku ora bisa \emph{mung} kandha yen
hierarki ngluwihi saben urutan becik. Prinsip
iku kudu nyatakake bab kang luwih kuwat.
Kanggo iki, \cite{Shoenfield:AST} ngusulake
penguatan kang lumrah: hierarki ora
\emph{kofinal} karo himpunan apa wae.\footnote{G\"odel
katon uga ngusulake gagasan kang mirip; delengen
\citet[p.~223]{Potter2004}. Kanggo rembugan
G\"odel lan \citeauthor{Shoenfield:AST},
delengen \citet[90--5]{Incurvati2020}.}
Luwih rinci, yen $\tau$ iku peta saka
himpunan menyang tataran hierarki, citra
himpunan apa wae $A$ miturut $\tau$ ora bakal
ngentekake kabeh hierarki. Kanthi skematis:
\begin{enumerate}
	\item[] \stagescofin. Yen $A$ himpunan lan
	$\tau(x)$ tataran kanggo saben $x \in A$,
	mula ana tataran kang luwih pungkasan
	tinimbang saben $\tau(x)$ kanggo $x \in A$.
\end{enumerate}
Cetha yen $\ZF$ mbuktekake versi
\stagescofin{} kang diformalake kanthi trep.
Kosokbaline, kanthi ora formal kita bisa
nerangake yen \stagescofin{} menehi dhasar
marang Panggantian.\footnote{Luwih angel
mbuktekake Panggantian saka sawijining
formalisasi \stagescofin, amarga $\Z$ dhewe
ora cukup kuwat kanggo netepake tataran-tataran;
mula ora cetha carane ngrumusake
\stagescofin kanthi formal. Salah siji pilihan
yaiku makarya
ing sawijining ekstensi saka $\LT$, kaya kang
dirembug ing
\olref[sth][replacement][extrinsic]{sec}.}
Upamane $(\forall x \in A)\lexists![y][\phi(x,y)]$.
Kanggo saben $x \in A$, tetepna $\sigma(x)$
minangka $y$ kang nyukupi $\phi(x,y)$, lan
$\tau(x)$ minangka tataran nalika $\sigma(x)$
pisanan kabentuk. Miturut \stagescofin, ana
tataran $V$ kang nyukupi
$(\forall x \in A)\tau(x)\in V$. Amarga saben
$\tau(x) \in V$ lan $\sigma(x) \subseteq \tau(x)$,
lan tataran iki poten, Pamisahan menehi
$\Setabs{y \in V}{(\exists x \in A)\sigma(x) = y}
= \Setabs{y}{(\exists x \in A)\phi(x,y)}$.

\begin{prob}
	Rumusna \stagescofin{} kanthi formal ing $\ZF$.
\end{prob}

Dadi \stagescofin{} ndhukung Panggantian.
Lan paling ora katon lumrah yen \stagesinex{}
ndhukung \stagescofin. Upamane \stagescofin{}
ora lumaku. Mula hierarki iku kofinal karo
sawijining himpunan~$A$: ana peta $\tau$
kang kanggo saben tataran~$S$ ana sawijining
$x \in A$ kang nyukupi $S \in \tau(x)$.
Yen mangkono, ana cara kanggo nyekel
``tanpa wates mutlak'' kang dianggep diduweni
hierarki: saben tataran \emph{kawengku} dening
sawijining anggota citra $\tau$ ing $A$.
Iki nglemesake gagasan
yen hierarki iku ``tanpa wates mutlak''.
Miturut kontraposisi, \stagesinex{}
ngakibatake \stagescofin, kang banjur menehi
dhasar marang Panggantian.

Iki upaya kang cukup njanjeni kanggo menehi
dhasar intrinsik marang Panggantian. Nanging
apa pungkasane bisa kasil, kita pasrahake
marang pamaca kanggo mutusake.

\end{document}
